\section{Introduction}
\label{sec:introduction}

Consider a binary segmentation system used with human oversight. Its mask can be
accepted as produced or sent for review, but the system must make that decision
before the reference annotation is available. This is especially consequential in
clinical workflows, where a plausible-looking contour may still miss a lesion
boundary or invent a distant component. Average test accuracy does not solve the
triage problem: two models with the same mean Dice can have very different value if
only one assigns low confidence to its costly failures. Selective prediction makes
this operational. A fixed model emits a mask \(\widehat Y(x)\) and an image-level
confidence \(g(x)\); accepting cases from high to low confidence produces a
risk--coverage curve.

The difficulty is that segmentation failure is not a single scalar notion. Dice and
IoU measure overlap, while HD95 measures boundary behavior. Moving a small component
far away may barely change Dice yet sharply increase surface error; a diffuse local
boundary error can reverse that comparison. Meanwhile, familiar scalar scores such
as mean maximum probability or entropy summarize pixelwise uncertainty without
specifying which harm they should rank. SDC supplies an important counterexample to
this metric-agnostic practice by aligning confidence with Dice. It also exposes the
next question: can metric alignment be extended beyond Dice without hiding the
probabilistic assumptions required by a nonlinear mask loss?

Our answer begins by separating three objects. The \emph{deployed action}
\(\widehat Y(x)\) is the mask that would actually be accepted; it may come from a
threshold of \(0.5\), \(0.3\), or any other fixed rule. The \emph{deployment loss}
\(L\) says what harm that action incurs. The remaining object is a law for the
unknown label mask, because conditional risk cannot be evaluated before annotation
without one. A probability map supplies pixelwise values, but it does not identify
their joint distribution.

We make this missing law explicit rather than implicit. Applying one shared threshold
\(T\sim\mathrm{Unif}(0,1)\) to the probability map generates a nested family
\(Y_T=\{i:p_i\geq T\}\). The resulting random set preserves \(p\) as its one-pixel
coverage function and converts expected geometric loss into a one-dimensional
threshold integral \citep{goodman,molchanov}. We call it a \emph{working posterior}:
it is transparent and computable, but is not claimed to be the true label posterior.
For quadrature nodes \(t_m\) and weights \(w_m\), it yields
\begin{equation}
  C_{L,M}(x)
  =-\sum_{m=1}^M w_m
    L\!\left(Y_{t_m},\widehat Y(x)\right).
  \label{eq:intro-score}
\end{equation}
Changing \(L\) now changes an overlap confidence into a boundary confidence while
leaving the deployed mask fixed. Every candidate label is compared with that same
action; dispersion among candidate labels would estimate something else. Under the
working posterior the integral target is exact, and midpoint quadrature controls its
numerical approximation. Crucially, increasing \(M\) cannot repair a miscalibrated
map or a wrong joint law. Our theory therefore separates deterministic quadrature
error from a loss-specific discrepancy between the working and true posteriors.

This separation leads directly to the empirical test. We hold each probability map
and deployed action fixed, construct Dice- and nHD95-indexed scores, and evaluate
both under matched and cross-loss AURC. Across ten model--dataset conditions,
nHD95-M32 outranks Dice-M32 for boundary risk in all ten, but their Dice-risk
ordering splits six to four; SDC remains the strongest primary Dice-risk score in
six conditions. The result supports loss indexing as a way to expose useful
geometric rankings, not as a promise that a same-named score always wins.

\subsection{Related Work}
\label{sec:related}

\paragraph{Selective prediction and failure detection.}
Selective classification formalizes the trade between coverage and risk
\citep{geifman2017selective,selectivenet,feng}, while failure-detection studies show
that conclusions depend on evaluation choices and score aggregation
\citep{jaeger,values}. We use AURC as their common ranking functional, but expose its
loss index because segmentation admits non-equivalent errors.

\paragraph{Segmentation confidence.}
ValUES systematizes how pixelwise uncertainty becomes image-level segmentation
confidence, and comparative benchmarking confirms that aggregation matters
\citep{values,zenk}. SDC is the closest metric-aligned predecessor: it uses soft
Dice between the probability map and deployed mask and derives a Dice-specific
approximation \citep{sdc}. We retain that deployed action, replace the
Dice-specific surrogate by loss-indexed expected risk, and state the joint working
posterior required by nonlinear mask losses.

\paragraph{Random sets and geometric losses.}
Random level sets have a long history in fuzzy and random-set theory
\citep{goodman,molchanov}, and shared-level constructions also arise in
probabilistic contour and excursion-set analysis \citep{poethkow,bolin}. We use
them as a transparent computational posterior, not an identification result.
Boundary-aware objectives and measures address behavior missed by overlap
\citep{karimi2020}; our experiment carries that distinction into confidence
ranking.

\subsection{Our Contributions}

Against this background, we first define loss-indexed selective risk \(\AURC_L\)
and its Bayes-optimal image-level confidence \(C_L^\star\), explicitly tying a
deployed action to its evaluation loss. We then introduce a single level-set
estimator for any bounded measurable segmentation loss. Its target is exact under
the stated working posterior, while midpoint quadrature has a deterministic
\(O(M^{-1})\) error bound separated from posterior-model error. For Dice, we show
why SDC and our score are distinct: SDC is a marginal-only ratio of expected
components, whereas level-set confidence is expected Dice under an explicit joint
law, and neither assumption set implies the other. Finally, we validate the
framework across ten fixed model--dataset conditions using tie-aware AURC and
paired, multiplicity-controlled inference. Together, these results make deployment
loss a principled axis for constructing and auditing confidence without claiming
that metric matching guarantees the best ranking.
